% Use, for example: \caption{\GrowthProfilesCaption}
% Protein captions are under paper/protein/figures/captions.tex.
\newcommand{\GrowthProfilesCaption}{%
Translation of the Gaussian $n(a,0)=100\exp[-(a-0.2)^2/(2(0.03)^2)]$
initially restricted to $0\le a\le1$, with $G=0.1$ and $t=4\,\mathrm{s}$.
Zero inflow and open outflow are used on 200 uniform bins. Panels (a,b) use
$\mathrm{CFL}=0.5$ and $\mathrm{CFL}=1$, respectively; the characteristic Inguva2022 update is
eligible only in (b). The initial centre moves from 0.2 to 0.6. Grey lines
show the analytical density and black points its exact bin averages.
Coloured lines join computed bin averages, with sparse markers at actual bin
centres; they do not add subcell resolution. The seven pairings match the
manual, including WENO5 with SSPRK3 and RK4 and conservative HOC with its
native midpoint update. Signed undershoots are retained without clipping.}

\newcommand{\GrowthWorkPrecisionCaption}{%
Smooth periodic translation on $[0,1]$, with
$n(a,0)=1+0.2\sin(2\pi a)+0.1\cos(4\pi a)$, $G=0.1$ and $t=4\,\mathrm{s}$.
Meshes have 40, 80, 160, 320 and 640 bins; $\mathrm{CFL}=0.5$ except for
Inguva2022's required $\mathrm{CFL}=1$. Initial and reference states are exact
bin integrals. (a) Integration wall time versus number of bins; (b) RMSE versus
the same measured wall time. RMSE is
$[\sum_i\Delta a_i(\bar n_i-\bar n_i^{\mathrm{exact}})^2/\sum_i\Delta a_i]^{1/2}$.
Each marker identifies one recorded mesh configuration. HOC uses its native
midpoint update, which is second order in time.
A separate compilation warmup precedes one measured solve per actual
method/boundary combination. Model/grid assembly, explicit GC, plotting and
file I/O are excluded; integrator initialization, stepping, checks, allocations
and saved states are included. There are no repeatability error bars.}

\newcommand{\PrecipitationFieldsCaption}{%
Transient number density in an axially dispersed continuous precipitator,
using Zhang et al. (2024), Case 7A, SI Table S8.
(a--d) $t=5,40,80,200\,\mathrm{s}$ in chronological order.
The computational logarithmic size edges span $1$--$900\,\mu\mathrm{m}$;
the linear size axis displays a $1$--$100\,\mu\mathrm{m}$ zoom.
The full size domain is retained for outlet curves, errors and inventories.
The axial mesh is uniform over $0$--$0.47\,\mathrm{m}$. The numerical reference
has 1000 size bins and 500 axial cells. The cell averages are
computed with geometry-aware size WENO5, axial MUSCL, and checked CVODE-BDF;
they are not CADET WENO35 results. CVODE uses adaptive BDF order up to five,
stability-limit detection and no prescribed restart,
relative tolerance $10^{-7}$, absolute tolerances $10^{-14}$ for the
$10^{12}$-scaled bin populations and $10^{-11}\,\mathrm{kg\,m^{-3}}$ for solute,
and a 0.5 s step cap. All panels share one linear color scale with the
\texttt{magma} palette, from zero to the maximum across the four full fields.
Zero and tiny negative tails use its darkest color; signed minima and negative
populations remain in the exports. No numerical states are clipped. No positive
startup background, image smoothing or spatial interpolation is introduced.
The negative-number fraction at the saved reference times is below $10^{-23}$.
This external integrator has no positivity guarantee; accepted states are checked.
The initial population and solute are zero. Primary and secondary nucleation,
solute depletion, Danckwerts inflow and zero-gradient outflow are retained.
Independent size counts 50, 100, 200 and 400 at fixed 500 axial cells, and axial
counts 50, 100, 200 and 400 at fixed 1000 size bins separate the two errors.
The same-mesh time-refinement check uses tolerances tenfold tighter and half the
step cap; its differences are exported at all four times. No analytical coupled PSD solution
exists for this case.}

\newcommand{\PrecipitationOutletCaption}{%
Outlet distributions at $t=200\,\mathrm{s}$ for the same Case 7A model.
(a) Absolute number density; (b) normalized volume-weighted size density,
$q_3(x_i)=x_i^3N_i/[\Delta x_i\sum_jx_j^3N_j]$, expressed per micrometre.
Here $x_i$ is the arithmetic bin centre.
The particle-size axis is logarithmic; densities remain per linear size,
not per log-size. Curves at axial counts 50, 100, 200 and 400 use 1000
logarithmic size bins, with a separately identified numerical reference solution.
The legend labels it ``Reference solution; it has $1000\times500$ cells.
Low axial resolution changes peak number and the distribution of solid material.
Open circles in (a), labelled ``Reference data, are digitized from the published Figure 9(c),
$N_x=N_z=200$, WENO35 in both coordinates; bars show conservative two-pixel
extraction uncertainty over the unambiguous 1.5--40 micrometre interval.
Published ordinate readings are unchanged and interpreted on the model's SI
m$^{-4}$ scale; taking the source figure's printed per-micrometre label
literally would exceed the feed's available material by orders of magnitude.
Only the number-density branch is digitized; its partial size window cannot
define an independently normalized volume distribution in (b).
The package uses size WENO5/axial MUSCL and the rounded SI flow rate
$1.67\times10^{-7}\,\mathrm{m^3\,s^{-1}}$; the public CADET input uses
$10\,\mathrm{mL\,min^{-1}}$. Volume densities integrate to one with actual
bin widths; they differ from the paper's material percentages per logarithmic
bin. Literature data and a finite numerical reference are not exact solutions.}

\newcommand{\PrecipitationRefinementCaption}{%
Separate refinement at $t=200\,\mathrm{s}$ against the common logarithmic-size
$1000\times500$ numerical reference. (a) Size counts 50, 100, 200 and 400 at
$N_z=500$; (b) axial counts 50, 100, 200 and 400 at $N_x=1000$.
Whole-field and outlet differences integrate the absolute difference of
piecewise-constant number densities over common partitions, divided by the
reference number integral on the corresponding domain. This differs from the
paper's spline comparison. Differences include representation, source-deposition,
evolution and finite-reference errors, not an isolated WENO order.
Both refinements reduce the outlet discrepancy; coarse axial grids give the
larger difference over these tested ranges. Near-overlapping squares and circles
are separately evaluated whole-field and outlet norms.}

\newcommand{\PrecipitationWorkPrecisionCaption}{%
The same two refinement paths and normalized differences as the preceding
figure, plotted against single-solve integration wall time. A separate small
compilation warmup precedes the measured solves for this method and boundary
combination. Stepping, checks and saved
states are timed; assembly, explicit GC, export and plotting are excluded.
Lines follow increasing resolution within a path and are not Pareto fronts.
There are no repeatability error bars.}
\newcommand{\MillingForwardCaption}{%
Constant-speed cone-mill startup compared with Barrasso et al. (2013),
DOI 10.1007/s12247-013-9155-0, Figure 7.
(a) Instantaneous volume-weighted outlet diameter median; (b) mill hold-up.
Solid curves use one conservative fixed-pivot breakage/classification scheme
and the checked CVODE-BDF solver at 3350 and 4923 RPM. There are 100 power-grid
volume pivots; the initial material is one second of feed (7.4 g), a declared
reconstruction assumption. The outlet median uses linear diameter interpolation
of the cumulative pivot-volume weights, the same convention used by the UQ
example. Open symbols, labelled ``Reference data'', are the unchanged published
simulation markers extracted from the original vector figure, not experimental
measurements. No kinetic parameters were fitted. The two speeds use the same
scheme; no additional mesh or observation-reconstruction curves are displayed.
This remains a partial reproduction: hold-up is close, while the diameter
trajectory differs. The UQ calculation below uses the same fixed-pivot/BDF
scheme with its separately declared geometric grid, feed deposition, empty
initial mill and speed change.}

\newcommand{\MillingUQCaption}{%
Forward uncertainty in the instantaneous volume-weighted outlet median of the
shared cone-mill model: (a) Monte Carlo, 12,000 realizations; (b) degree-one
Hermite PCE, 12 LHS training inputs; (c) degree-two Hermite PCE, 32 hybrid training
inputs. Four Gaussian kinetic/classification inputs are assumed independent.
Solid curves are propagated means; grey dashed curves use the parameter means.
Shading gives pointwise 2.5/97.5 percentiles; dotted curves give the explicitly
approximate mean $\pm1.96$ output SD. These describe output uncertainty, not
confidence intervals on a mean or simultaneous trajectory coverage. PCE
percentiles use 50,000 seeded surrogate evaluations under the full Gaussian
law; bounded training designs do not truncate that law. All panels use the same
axes and retained numerical summaries as the documentation case study.
The mill starts empty on 100 geometric volume pivots from
$6.54\times10^{-17}$ to $10^{-3}\,\mathrm{m^3}$; the first displayed time is
0.001 s because the empty-start outlet median is undefined at zero.
Speed changes from 4923 to 1500 RPM at 30 s, resolved by two BDF integrations.
The conservative model uses relative tolerance $10^{-9}$ and absolute tolerance
$10^{-14}$ on per-pivot/discharged material in kg. The Bhonsale et al. (2018)
DOI 10.3390/pr6120255 design/band ambiguities and declared initialization,
domain and observation choices prevent a claim of exact published-curve
reproduction. Original ensembles, producer records and timings are reused;
this redraw performs no new MC or timing measurements.}
